\section{Fixed continuous uncertainty and the price of transport evidence}\label{app:probability}
A single unknown physical parameter can influence every conditional transition. A finite list of mixture weights need not describe this situation. We first ask which controller guarantees remain decidable, then what an independent check of only \emph{part} of a row can establish.

\subsection{Decision with a fixed physical parameter}
A semialgebraic set is described by finite Boolean combinations of polynomial equalities and sign inequalities with rational coefficients. A world consists of a finite regime and a fixed vector \(\theta\) in its declared semialgebraic region. Finite-horizon initial and conditional laws, labels and bounded payoffs have semialgebraic graphs; any denominator is certified positive on its region. The complete normalized conditional rows, real event flags, world membership, implementation and support rules of Theorem~\ref{flag:p1} apply. Private randomness precedes no unauthorized disclosure. Separate the exact guard class \(W\) from the class \(C\) used for performance requirements. Let \(X\) be the compact sequence polytope of that theorem and \(d=\dim(\operatorname{aff}X)\) if nonempty.

\begin{theorem}[G01.1. Algebraic controllers and finite obstructions]\label{adv:g011}
The robust problem has an exact decision procedure. For a feasible system it returns real algebraic realization factors with certificates that no admissible world violates a hard or soft constraint. For infeasibility with \(X\ne\varnothing\), at most \(d+1\) algebraic worlds, with their applicable linear constraints, suffice to witness infeasibility. It also decides label homogeneity, bounded score extrema and their attainment for a specified algebraic controller. A bounded optimal worst cost or utility on a nonempty robust policy set is attained with algebraic policy and value, even if the worst world is only approached.
\end{theorem}
\begin{proof}
The chance-weight product at a fixed leaf is semialgebraic in the same \(\theta\) throughout. Multiply by that leaf's sequence factor and sum. Scores are semialgebraic jointly and affine in the controller at each world. A nonnegative invalid marker of expectation zero on \(W\) enforces the whole support, including isolated parameter points. Feasibility has the real-field form \(\exists x\in X\,\forall w\,[\mathrm{hard}(x,w)\wedge\mathrm{soft}(x,w)]\), restricted to the applicable classes. Quantifier elimination decides this formula and gives a real algebraic sample in each nonempty decision cell. Encode \(f\ge0\) as \(f=t^2\), and \(f>0\) as \(ft^2=1\), after splitting Boolean cases and introducing graph variables. For an empty system of equalities \(f_j=0\), the real Nullstellensatz supplies an identity \(-1=\sum_i q_i^2+\sum_j u_j f_j\) over real algebraic coefficients. Exact root isolations order those coefficients.

The robust set is an intersection of closed halfspaces in compact \(X\). Infeasibility yields a finite subintersection by compactness and \(d+1\) worlds by Helly. Search for that many admissible worlds and a linear infeasibility multiplier by the same decision procedure. The retained witnesses are algebraic, physically meaningful as countermodels only with realization. Bounded worst cost is lower semicontinuous in \(x\), since it is a supremum of continuous affine functions. It attains a minimum on the robust set. Quantifying bounds, arbitrarily close witnesses, and attainment gives its exact algebraic value; reverse signs for utility. Compact continuous world domains attain inner extrema. Compactness of the world domain alone does not imply this for discontinuous rows. The computational cost can be large, and a physical sampler for an algebraic lottery requires its own resource model.
\end{proof}

\begin{theorem}[G01.2. An approximation with a checked margin]\label{adv:g012}
Keep the exact hard guard on all its worlds. Let every soft score vary by at most \(Lr\) between any world and an \emph{actual} mesh member in its regime, uniformly over \(X\). Imposing thresholds \(b_i-Lr\) on the mesh suffices for exact feasibility; thresholds \(b_i\) there are necessary. If one hard-legal policy has uniform slack \(s>0\), mixing it into a relaxed mesh solution with weight \(Lr/(s+Lr)\) yields a feasible policy. Linear cost of range \(D\) increases by at most \(DLr/(s+Lr)\).
\end{theorem}
\begin{proof}
The deviation inequality extends a tightened check to every world and bounds the relaxed policy's possible excess by \(Lr\). Affine mixing gives excess at most \((1-t)Lr-ts=0\) at the stated weight \(t\). Both component policies obey the exact hard guard. Bounding the cost by its range proves the claim.
\end{proof}
\paragraph{Why continuous inference needs care.}
The constraints \(x\le\theta,\ x\ge1/(2\theta)\) on all \(1/\sqrt2\le\theta\le1\) force the unique policy \(x=1/\sqrt2\) with rational defining data. The constraint \(\theta x\le\theta^2\) for every \(0\le\theta\le1\) forces \(x=0\), yet any finite selection allows positive \(x\); restricting \(x\) to \((0,1]\) even destroys the finite infeasibility conclusion. Two independent draws conditional on the same \(\theta\) have first and joint moments \((\theta,\theta^2)\). Every point of that parabola is exposed in its convex hull by its own tangent. Hence no finite polytope of complete laws preserves all linear payoff envelopes. For a graph with \(m\) edges choose a uniformly random edge \(ij\) and score disagreement of independent draws with means \(\theta_i,\theta_j\). The expected score is \(\sum_{ij}(\theta_i+\theta_j-2\theta_i\theta_j)/m\). It is affine in each coordinate separately, hence maximized at a cube vertex and equal there to the maximum cut fraction. A universal upper-threshold test therefore includes the complement of the maximum-cut decision problem. This is coNP-hard even for a single policy at constant depth.

\subsection{Comparing the parts of two laws that evidence actually covers}
A subprobability measure is nonnegative mass with total at most one. Fix nominal and deployment sources with a common legal controller and matched pairs of checkpoints whose complete received records agree. On each side all supported paths, including unmatched paths, meet the physical permission and hard-budget requirements. At the start and then at every matched history, action and world, a verified joint subrow on paired successors has each marginal bounded by its side's genuine conditional row. The matched pair is chosen before the next random outcome. Keep the same installed world and independent policy seed throughout.

\begin{theorem}[G02.1. Matched mass and a selective score]\label{adv:g021}
Multiply the common subrows and common decision factors down the tree to obtain terminal pair mass \(\mu\) of size \(m\). For a genuine score in a common range \([L,U]\), define \(s=\sum\mu(F_D-F_N)\) on the matched terminals, where subscripts denote deployment and nominal. The difference of true expected scores is within \((U-L)(1-m)\) of \(s\). If only a deployment score interval \([\ell_t,u_t]\) is known on matched terminal \(t\), its true expectation is in
\[
 \left[\sum_t\mu(t)\ell_t+L(1-m),\
       \sum_t\mu(t)u_t+U(1-m)\right].
\]
If every reached subrow at level \(j\), including the initial row, retains a fraction at least \(a_j\in[0,1]\), then \(m\ge\prod_j a_j\).
\end{theorem}
\begin{proof}
Induct on prefix length. Equal records allow the same action probabilities. Conditional marginal domination gives domination of both next genuine prefix masses. Each side thus has residual mass \(1-m\). Its possible contribution lies between \(L(1-m)\) and \(U(1-m)\). Subtracting the two residual bounds proves the discrepancy; multiplying the lower survival bounds proves the product statement. Sharpness requires compatible residuals attaining the endpoints. Truncate a vacuous negative survival lower bound to zero. For \(0\le\varepsilon,u\le1\), the scores \(I(V-\varepsilon)\) and \(I(G-u)\) have width one, including value zero off issuance. Upper risk score zero, lower utility score zero and \(\mathbb EI>0\) imply the conditional guarantees. Use the actual score range for other normalizations.
\end{proof}

\begin{theorem}[G02.2. Which source comparisons must be bought?]\label{adv:g022}
Let a finite exact read catalog have explicitly listed full report tuples, semialgebraic world regions and nonnegative prices. Reads are safe, once-only, nondisturbing. Each received report may filter worlds, add a matched-mass fragment or tighten a matched payoff interval. The combined fragments must be jointly dominated on both sides. Keep the exact hard guard fixed wherever its physical premises require the original guard. Positive certification requires label homogeneity and one common policy satisfying the G02.1 bounds; negative certification requires homogeneous negativity and lawful nonactivation. A finite counterexample-guided minimum hitting-set procedure finds the least sufficient realized price in this proof family. Testing all report tuples gives the minimum universal fixed set. A minimax recursion over partial reports gives the minimum prospective worst price.
\end{theorem}
\begin{proof}
Before filtering numerical worlds, compile any invariant hard guard into the policy domain. A positive coefficient on an invalid leaf in some guarded world forces its nonnegative sequence factor to zero. Equivalently retain a verified fixed hard subpolytope in each witness test. At a failed subset \(S\), obtain an opposite world or at most \(d+1\) retained worlds blocking the affine envelope tests by Theorem~\ref{adv:g011}. An unread field can help such a group by excluding a member, adding positive matched mass somewhere on it or tightening an interval endpoint there. Collect all such indices in a cut. Any further set omitting every one of them, even after union with \(S\), leaves the blocker and its failed tests; a sufficient set must hit the cut. Extra independently dominated fragments and tightened score intervals can only improve the admissible enclosing bounds. An empty cut proves catalog insufficiency. A minimum hitting-set proposal followed by an exact separation oracle yields an attaining certificate at termination; each failure adds a previously unhit subset of the finite menu. The same argument across tuples yields the fixed list, while conditioning on every possible next reply gives the prospective recursion. Numeric narrowing never authorizes a command that was forbidden somewhere in the retained hard guard.
\end{proof}

\begin{theorem}[G02.3. Safe adaptively selected diagnostics]\label{adv:g023}
For each registered diagnostic \(j=1,\ldots,k\), suppose actual observations lie in \([0,1]\) and their conditional means at every selected pre-draw cut differ from a specified \(m_j(w)\) by at most \(b_j\). For \(0<\eta<1\) choose a rational outward radius \(r_n\) between \(\sqrt{\log(2k n(n+1)/\eta)/(2n)}\) and that number plus \(1/n\). At each reached count \(n\), intersect the class with the interval of radius \(b_j+r_n\) around the observed mean, clipped to \([0,1]\). The true world survives simultaneously with probability at least \(1-\eta\) under any predictable selection and adapted stop. If neighborhoods of the actual moments with radii \(\lambda_j>2b_j\) all allow the required common decision, quotas with \(2(b_j+r_{n_j})\le\lambda_j\) attain success probability at least \(1-\eta\), provided those quotas are lawful.
\end{theorem}
\begin{proof}
For the reached observations of one stream let \(S\) sum the deviations from their actual conditional means and \(N\) count them. The conditional bounded-draw moment inequality makes \(\exp(tS-t^2N/8)\) a nonnegative supermartingale for all real \(t\). To see the inequality, bound the conditional exponential by that of a Bernoulli with the same mean; its centered log moment has zero initial slope and second derivative at most \(1/4\). Stopping the exponential at first crossing and any finite time proves the crossing bound by expectation. Thus the probability of reaching count \(n\) with either tail at least \(nr\) is at most \(2e^{-2nr^2}\), taking \(t=\pm4r\). Sum \(\eta/[k n(n+1)]\) over streams and counts. On inclusion, all retained moments differ from the true moments by at most twice the radius including bias. This yields the quota claim under the unchanged guard. If a separate label mechanism can fail with probability \(\delta\), the mass of issued runs subject to either failure is at most \(e=\min(a,\eta+\delta)\), where \(a=\Pr(I)\). A fresh correctly linked deployment then gives \(\Pr(V\cap I)\le\varepsilon a+(1-\varepsilon)e\) and \(\mathbb E[IG]\ge u(a-e)\). Conditional interpretations require \(a>0\). An unverified instrument or model membership is outside this probability statement.
\end{proof}

\begin{proposition}[G02.4. A nonzero residual with sharp ratios]\label{adv:g024}
A five-tag source has masses \(3/5,1/5,1/20,1/20,1/10\). On the first two tags a verified positive report makes the work a harmless useful activation; a zero report suppresses it. The next two tags respectively activate with a completed but violating job and do not activate. In the last tag one Bernoulli draw activates with chance \(\theta\), and a second, conditionally independent given this same \(\theta\in[0,1]\), causes a violation with chance \(\theta\); it earns no utility. Two nonchanging source reads of prices one and two verify the routine tags. The intrinsic target is positive exactly when both reports are one. At a positive report least realized setup cost is three, and the best fixed and prospective worst costs are three. The always-start policy gives activation at least \(17/20\), selective risk at most \(3/19\), and conditional utility at least \(17/19\).
\end{proposition}
\begin{proof}
The two checks plus the known middle tags certify mass \(9/10\); known activation, joint violation and useful activation are \(17/20,1/20,17/20\). On the full actual source those quantities are \((17+2\theta)/20,(1+2\theta^2)/20,17/20\). At thresholds \(3/19\) and \(17/19\) the upper risk score and lower utility score from Theorem~\ref{adv:g021} both equal zero, with equality in both actual ratios at \(\theta=1\). Dividing the separate worst numerator by the separate lowest denominator gives the looser \(3/17\). Missing either routine tag allows loss of at least \(1/5\) of activation, which even a full \(1/10\) residual cannot restore at the floor. Thus any successful positive path must buy both reads, and buying both suffices everywhere, with early negative stops available. No measurement of the two unchecked conditional probabilities is needed.
\end{proof}
